\section{Il Ruolo dei Dataset}
\label{sec:dataset}

Lo sviluppo e il progresso di qualsiasi campo del \textit{deep learning} sono strettamente legati alla disponibilità di \textbf{dataset} di alta qualità, diversificati e rappresentativi della complessità degli scenari reali. I dataset non costituiscono semplici collezioni di immagini o campioni, ma rappresentano la base per definire problemi, formulare ipotesi scientifiche, addestrare e validare gli algoritmi, misurare i progressi ottenuti. L'efficiacia dei metodi può essere valutata solo grazie a benchmark condivisi e rigorosi. 

Nel contesto specifico dell'\textbf{anomaly detection} per applicazioni industriali, i dataset assumono un'importanza ancora più critica. A differenza di altri domini dove i dati possono essere raccolti in modo relativamente scalabile (ad esempio, attraverso il \textit{web scraping} o l'acquisizione da sensori diffusi), le immagini di prodotti industriali, specialmente quelle che riportano difetti difficili da riprodurre, richiedono acquisizione controllata e spesso condotta in am\-bien\-ti produttivi reali. Inoltre, la natura stessa dei difetti industriali (variabili da micro-graffi superficiali a difetti strutturali complessi, da anomalie di texture ad anomalie logiche) impone che i dataset siano progettati con cura per coprire uno spettro ampio e realistico di possibili scenari.

% ci sta?
Un dataset di qualità per l'anomaly detection industriale deve soddisfare alcuni requisiti fondamentali \cite{bergmann2019mvtec, bergmann2022beyond}:

\begin{itemize}
\item \textbf{Immagini ad alta risoluzione}: La rilevazione di micro-difetti (graffi sottili, piccole contaminazioni) richiede immagini con sufficiente dettaglio spaziale per preservare le caratteristiche distintive delle anomalie.

\item \textbf{Annotazioni pixel-precise}: Per applicazioni industriali che richiedono localizzazione accurata dei difetti, sono necessarie maschere di \textit{ground truth} con segmentazioni precise dei confini delle anomalie.

\item \textbf{Variabilità controllata}: I dati normali devono presentare variabilità intra-classe sufficiente (variazioni di illuminazione, pose, texture) per addestrare modelli robusti, senza compromettere la definizione di normalità.

\item \textbf{Diversità di anomalie}: Le tipologie di difetti devono coprire un ampio spettro di manifestazioni visive: difetti strutturali (crepe, denti), difetti superficiali (graffi, macchie), anomalie logiche (componenti mancanti o disallineati).

\item \textbf{Scalabilità}: I dataset devono essere sufficientemente grandi per permettere l'addestramento di modelli \emph{deep learning} e una validazione statistica robusta dei risultati.
\end{itemize}

%

L'evoluzione dei metodi di anomaly detection, illustrata in \ref{sec:metodi}, è andata di pari passo con l'evoluzione dei dataset a disposizione. I primi approcci, spesso basati su \textbf{autoencoder} o \textbf{GAN}, sono stati valutati su dataset relativamente piccoli e omogenei (ad esempio, immagini in scala di grigi di singole componenti). La pubblicazione di dataset come \textbf{MVTec AD} ha segnato una svolta, offrendo una varietà di oggetti e texture, annotazioni pixel-precise e una chiara separazione tra training (solo esempi normali) e test (con anomalie). Questo ha permesso la nascita e il confronto di metodi più sofisticati, basati su \textit{feature pre-addestrate}, \textit{memory bank} e \textit{normalizing flows}. 


Più recentemente, la necessità di affrontare scenari industriali sempre più complessi (ispezioni 3D, illuminazione variabile, oggetti multi-vista) ha portato verso la creazione di dataset \textbf{multimodali} (RGB, depth, point cloud) e \textbf{large-scale}, spesso acquisiti in ambienti reali o generati sinteticamente con alto realismo. Questi dataset consentono di testare la robustezza e la generalizzazione dei metodi esistenti, ma portano anche verso lo sviluppo di nuove architetture in grado di elaborare informazioni visive e geometriche in modo integrato.

In questo capitolo, presenteremo una panoramica strutturata dei principali dataset utilizzati nella letteratura di \emph{anomaly detection} industriale. Dopo averne analizzato le caratteristiche distintive e le metriche di valutazione associate, discuteremo come l'evoluzione dei dataset abbia influenzato lo sviluppo di algoritmi sempre più accurati, robusti e pronti per il deployment in scenari industriali reali.

%% 3.2 Evoluzione
\section{Evoluzione dei Dataset (2018–2025)}
\label{sec:evoluzione_dataset}

L'evoluzione dei dataset per l'anomaly detection industriale riflette la crescente complessità delle applicazioni industriali e l'emergere di nuove esigenze. Possiamo identificare tre fasi principali, ciascuna caratterizzata da un salto in termini di modalità di acquisizione, tipologia di anomalie e scenario valutativo.

\subsection{Prima generazione (2018–2020): Benchmark 2D e definizione del task.}
I primi dataset, come \textbf{Magnetic Tile Defects} (2018) \cite{huang2018surface} e \textbf{MVTec AD} (2019) \cite{bergmann2019mvtec} si concentrarono su immagini 2D (scala di grigi o RGB) di oggetti isolati o texture, con annotazioni pixel-precise per difetti superficiali (graffi, crepe, macchie). MVTec AD in particolare stabilì lo standard per dataset 2D RGB con annotazioni pixel-level, diventando il benchmark di riferimento universale.


\subsection{Seconda generazione (2021–2022): Espansione mul\-ti\-mo\-da\-le e anomalie complesse.}
Per testare modelli più avanzati e affrontare limiti dei dati puramente visivi, emersero dataset che incorporano informazioni 3D e anomalie di tipo logico\-strutturale. \textbf{MVTec 3D-AD} (2021) \cite{bergmann2021mvtec3d} introdusse depth map e point cloud insieme alle immagini RGB. \textbf{MVTec LOCO AD} (2022) \cite{bergmann2022beyond} focalizzò l'at\-ten\-zio\-ne sulle \textit{logical anomalies} (componenti mancanti, assemblaggi errati), ri\-chie\-den\-do ai modelli di ragionare sulla struttura globale dell'oggetto. Parallelamente, dataset sintetici come \textbf{Eyecandies} (2022) \cite{bonfiglioli2022eyecandies} permisero un controllo fine su illuminazione, materiali e geometrie, facilitando la validazione in condizioni perfettamente controllate.

\subsection{Terza generazione (2023–2025): Scenari realistici, multi-vista e large-scale.}
I dataset più recenti hanno spostao l'attenzione verso condizioni realistiche. Dataset come \textbf{Real-IAD} (2024) \cite{ding2024realiad} e \textbf{RAD} (2024)  \cite{liu2024rad} offrono migliaia di immagini acquisite da più viewpoint, con oggetti riflettenti, trasparenti e in condizioni di illuminazione variabile. \textbf{MVTec AD 2} (2025) \cite{bergmann2025mvtecad2} introduce sfide aggiuntive come difetti ai bordi dell'immagine e variazione nel numero di oggetti per scena. L'obiettivo è ridurre il gap tra performance su benchmark controllati e deployment industriale reale.

La Tabella \ref{tab:dataset_evolution} sintetizza questa progressione temporale, evidenziando l'in\-cre\-men\-to di complessità e realismo.

\begin{table}[htbp]
\centering
\caption{Evoluzione dei dataset per anomaly detection industriale (2018-2025). La suddivisione in tre generazioni riflette l'espansione dalle basi 2D alla multimodalità, fino agli scenari realistici e multi-vista.}
\label{tab:dataset_evolution}
\footnotesize
\begin{tabular}{>{\raggedright\arraybackslash}p{0.8cm}>{\raggedright\arraybackslash}p{2.5cm}>{\raggedright\arraybackslash}p{10cm}}
\toprule
\textbf{Anno} & \textbf{Dataset} & \textbf{Contributi Principali / Innovazioni} \\
\midrule
\multicolumn{3}{c}{\textit{\textbf{Prima Generazione: Benchmark 2D e definizione del task (2018--2020)}}} \\
\midrule
2018 & Magnetic Tile Defects & Dataset pionieristico per difetti industriali; immagini in scala di grigi; valutazione di modelli di saliency detection. \\
\addlinespace
2019 & \textbf{MVTec AD} & \textbf{Benchmark di riferimento universale}; 15 categorie (oggetti e texture); annotazioni pixel-precise; separazione chiara training/test; standardizzazione di metriche (AUROC, PRO). \\
\addlinespace
2021 & MPDD & Focus su parti metalliche; scenari tipici dell'industria metallurgica (graffi, fori, ruggine). \\
\addlinespace
2022 & VisA & Oggetti con strutture complesse e multiple istanze; introduce protocolli few-shot e high-shot per valutare la sensibilità locale dei metodi self-supervised. \\
\midrule
\multicolumn{3}{c}{\textit{\textbf{Seconda Generazione: Espansione multimodale e anomalie complesse (2021--2022)}}} \\
\midrule
2021 & MVTec 3D-AD & Primo dataset 3D industriale; unisce RGB, depth map e point cloud; anomalie geometriche (graffi, ammaccature, contaminazioni). \\
\addlinespace
2022 & MVTec LOCO AD & Introduzione delle \textbf{anomalie logiche} (errori di assemblaggio, componenti fuori posto); metrica sPRO per valutare il rilevamento di difetti contestuali. \\
\addlinespace
2022 & Eyecandies & Dataset \textbf{sintetico} su larga scala (90K campioni); controllo completo su illuminazione, materiali e anomalie; include RGB, depth e mappe delle normali. \\
\midrule
\multicolumn{3}{c}{\textit{\textbf{Terza Generazione: Scenari realistici, multi-vista e large-scale (2023--2025)}}} \\
\midrule
2023 & Real3D-AD & Dataset 3D ad alta precisione; setting \textbf{few-shot} (4 campioni per categoria); copertura a 360° senza punti ciechi. \\
\addlinespace
2023 & PAD / MAD & Setting \textbf{pose-agnostic}; acquisizioni multi-vista ($\sim$20 viewpoint); versioni sia reali (MAD-Real) che sintetiche (MAD-Sim). \\
\addlinespace
2024 & Real-IAD & Dataset \textbf{large-scale} (150K immagini); 30 categorie industriali; acquisizione multi-vista (5 angolazioni); progettato per il deployment in scenari reali. \\
\addlinespace
2024 & RAD & \textbf{Multi-vista estremo} (60+ viewpoint); setting \textbf{noisy-pose} per testare la robustezza; include ricostruzione 3D a partire dalle viste RGB. \\
\addlinespace
2025 & \textbf{MVTec AD 2} & Introduce \textbf{sfide realistiche}: oggetti trasparenti/riflettenti, condizioni di illuminazione variabili, anomalie ai bordi dell'immagine, variazione nel numero di oggetti. \\
\addlinespace
2025 & Real-IAD D³ & Estensione \textbf{tri-modale} di Real-IAD: RGB, point cloud e mappe delle normali; 20 categorie; 69 gruppi di difetti per valutazione fine. \\
\addlinespace
2025 & SiM3D & Dataset \textbf{ibrido} (reale + sintetico); multi-vista (12--36 viste per istanza); valutazione voxel-level per segmentazione 3D. \\
\bottomrule
\end{tabular}
\end{table}




%% section metriche
\subsection{Metriche di Valutazione: Panoramica}
\label{subsec:metriche_intro}

La valutazione rigorosa degli algoritmi di anomaly detection richiede metriche standardizzate che catturino aspetti diversi della performance: accuratezza a livello di immagine, precisione nella localizzazione, robustezza ai falsi positivi, etc. Le metriche più utilizzate nei benchmark industriali possono essere classificate in due categorie principali: \textit{threshold-free} (indipendenti da una soglia di decisione) e \textit{threshold-dependent} (che richiedono la binarizzazione dello score di anomalia).

\paragraph{Metriche Threshold-Free.}
\begin{itemize}
\item \textbf{AUROC (Area Under the Receiver Operating Characteristic Curve)}: misura la capacità del modello di separare classi normali e anomale al variare della soglia. Viene riportata sia a livello immagine (\textit{Image AUROC}) che a livello pixel (\textit{Pixel AUROC}). Valori prossimi a 1 indicano prestazioni eccellenti. È la metrica più comunemente riportata per la sua robustezza a classi sbilanciate.
\item \textbf{PRO (Per-Region Overlap)} \cite{bergmann2022beyond}: progettata per valutare la qualità della localizzazione, specialmente per difetti piccoli. Calcola la curva Precision-Recall media su ogni regione anomala indipendentemente dalla sua dimensione, evitando che difetti grandi dominino la metrica. Il suo integrale (AUPRO) è diventato standard per dataset con annotazioni pixel-precise.
\end{itemize}

\paragraph{Metriche Threshold-Dependent.}
\begin{itemize}
\item \textbf{F1-Score}: media armonica di precision e recall, calcolata dopo aver applicato una soglia ottimale (tipicamente sul validation set). Sensibile allo sbilanciamento delle classi, ma utile per scenari dove è critico bilanciare falsi positivi e falsi negativi.
\item \textbf{MAE (Mean Absolute Error)}: utilizzata principalmente in metodi basati su ricostruzione o saliency map, misura l'errore medio assoluto tra la mappa predetta e la ground truth.
\end{itemize}

La scelta della metrica influenza direttamente la comparabilità dei risultati. Dataset incentrati sulla localizzazione (es. MVTec AD) privilegiano AUROC e PRO, mentre dataset con anomalie logiche (es. MVTec LOCO AD) possono introdurre metriche ad-hoc come \textit{sPRO} (Saturated PRO). Nei paragrafi seguenti, per ciascun dataset indicheremo le metriche con cui è stato principalmente valutato, collegandole alle specifiche esigenze applicative che lo hanno ispirato.

%%% DATASET 2D
\section{Dataset 2D RGB}
I dataset basati su immagini RGB 2D rappresentano il nucleo della ricerca in anomaly detection visiva. Questa sezione presenta i principali benchmark ordinati cronologicamente, evidenziando le caratteristiche distintive e i contributi di ciascuno.

% surface detection
\subsection{Surface Defect Saliency of Magnetic Tile (2018)}
\label{sec:magnetic_tile}
Rappresentante della prima generazione di dataset, Magnetic Tile \cite{huang2018surface} è stato uno dei primi dataset specificamente creati per anomaly detection industriale, focalizzato sul rilevamento di difetti superficiali in piastrelle magnetiche.

\paragraph{Caratteristiche.} Il dataset contiene \textbf{1,344 immagini in scala di grigi} di piastrelle magnetiche distribuite in 6 classi a seconda del tipo di difetto: Blowhole, Crack, Fray, Break, Uneven, e Free (assenza di difetti). Le immagini catturano difetti tipici del processo di produzione di componenti magnetici.

\paragraph{Contributo Metodologico.} Il dataset introdusse \textbf{MCuePushU}, un mo\-del\-lo di saliency detection che combina:
\begin{enumerate}
\item Estrazione di \textit{dominant cues} (caratteristiche salienti) dai dati
\item Fusione di questi cues in una rete U-Net \cite{ronneberger2015u} tramite operazioni aritmetiche sulle immagini
\item Un modulo Push Network che produce bounding boxes per evidenziare i difetti predetti
\end{enumerate}

\paragraph{Valutazione.} Gli esperimenti confrontarono MCuePushU con i principali me\-to\-di non deep-learning dell'epoca su dataset di scene naturali: GMR \cite{yang2013saliency}, MBP \cite{zhang2013minimum}, MC \cite{jiang2013automatic}, RBD \cite{zhu2014saliency}, DRFI \cite{jiang2013salient}. Le metriche utilizzate furono:
\begin{itemize}
\item \textbf{AUC (Area Under ROC)}: misura la capacità di separare regioni normali e anomale
\item \textbf{MAE (Mean Absolute Error)}: errore medio assoluto tra saliency map e ground-truth
\item \textbf{F$_\beta$}: massimo della curva Precision-Recall, con $\beta$ che bilancia precision e recall.
\end{itemize}

\paragraph{Limitazioni.} Sebbene pioneristico, il dataset presentava limitazioni significative: immagini in scala di grigi (perdita di informazioni cromatiche), dimensione limitata ($<1,500$ immagini), focus su un singolo tipo di oggetto. Tuttavia, il suo contributo fu fondamentale nel dimostrare la necessità di dataset industriali dedicati e nel stabilire protocolli di valutazione che sarebbero stati adottati e ampliati da benchmark successivi come MVTec AD.

% mvtec
\subsection{MVTec AD (2019)}
\label{sec:mvtec_ad}
MVTec Anomaly Detection (MVTec AD) \cite{bergmann2019mvtec} è diventato il benchmark di riferimento universale per l'anomaly detection industriale, utilizzato come baseline in praticamente tutti i lavori successivi.

\paragraph{Caratteristiche.} Il dataset contiene \textbf{5,354 immagini RGB ad alta ri\-so\-lu\-zio\-ne} distribuite su 15 categorie distinte: 10 categorie di oggetti (bottle, cable, capsule, hazelnut, metal nut, pill, screw, toothbrush, transistor, zipper) e 5 categorie di texture (carpet, grid, leather, tile, wood). Le immagini normali (3629 defect-free) costituiscono il training set, mentre il test set include 1725 immagini con anomalie. Sono presenti \textbf{73 tipologie di difetti} generati manualmente (in media 5 per categoria), riproducendo difetti realistici: graffi, denti, contaminazioni, deformazioni, componenti mancanti, disallineamenti

\paragraph{Annotazioni Pixel-Precise.} Ogni immagine anomala è corredata da una \textit{ground truth} mask binaria che segmenta a livello di pixel la regione difettosa. Questo permette la valutazione sia del task di \textbf{image-level classification} (l'im\-ma\-gi\-ne è normale o anomala?) che di \textbf{pixel-level localization} (dove si trova esattamente il difetto?).

\paragraph{Protocollo Sperimentale.} 
\begin{itemize}
    \item \textbf{Training}: solo su immagini normali (semi-supervised/one-class)
    \item \textbf{Testing}: su immagini normali e anomale, con valutazione separata per detection e localization
    \item \textbf{Data augmentation}: random crop di patch rettangolari ruotate per le texture, mentre per gli oggetti si applicano traslazioni, rotazioni casuali e mirroring dove possibile, generando 10000 patch di training per categoria.
\end{itemize}

\paragraph{Metodi Valutati e Performance Iniziali.} Il paper originale valutò diversi approcci rappresentativi dell'epoca, rivelando limiti significativi:
\begin{itemize}
    \item \textbf{AnoGAN} \cite{schlegl2017unsupervised}: GAN-based, discusso in Sezione \ref{subsec:metodi-GAN}. 70\% image AUROC, in difficoltà nelle categorie che presentano oggetti con più variazioni (forma, orientamento).
    \item \textbf{Autoencoder} (MSE/SSIM): varianti con loss L2 (MSE) e SSIM (Structural Similarity), discussi in Sezione \ref{subsec:autoencoder}. 80-85\% image AUROC (problemi di false positive, falliscono nel modellare piccoli dettagli).
    \item \textbf{CNN Feature Dictionary}: metodo basato su feature pre-addestrate, performance migliori ma limitate
    \item \textbf{Texture Inspection GMM-Based}: efficace per texture ma non per oggetti
    \item \textbf{Variation Model}: adatto a oggetti con variazioni limitate
\end{itemize}

\paragraph{Metriche Standardizzate.} Il dataset introdusse e popolarizzò due metriche fondamentali:
\begin{itemize}
    \item \textbf{AUROC (Area Under ROC Curve)}: per valutazione a livello im\-ma\-gi\-ne e pixel
    \item \textbf{PRO (Per-Region Overlap)}: metrica pixel-level che penalizza errori su piccole regioni anomale, affrontando il bias dell'AUROC pixel-level verso grandi regioni (discussa in Sezione \ref{sec:metriche})

\end{itemize}

\paragraph{Impatto.} MVTec AD ha avuto un profondo impatto sulla ricerca:
\begin{enumerate}
    \item Definì uno standard riproducibile e condiviso per la valutazione
    \item Espose i limiti dei metodi contemporanei (GAN-based, autoencoder)
    \item Motivò lo sviluppo di approcci più avanzati che oggi superano il 99\% AUROC
    \item Divenne la piattaforma di riferimento per confronti oggettivi e riproducibili
\end{enumerate}

\paragraph{Disponibilità e Adozione.} Distribuito sotto licenza Creative Commons, MVTec AD è ampiamente utilizzato sia in ricerca che in contesti industriali, con implementazioni integrate nei principali framework di deep learning e continui aggiornamenti (come MVTec AD 2, 2025).

% mpdd
\subsection{MPDD (2021)}
\label{sec:mpdd}

Metal Parts Defect Dataset (MPDD) \cite{jezek2021deep} fu introdotto per colmare una lacuna specifica: la scarsità di dataset focalizzati su parti metalliche, un dominio con caratteristiche distintive (superfici riflettenti, texture uniformi, variazioni cromatiche sottili). 

\paragraph{Caratteristiche.} MPDD contiene \textbf{1,346 immagini RGB} distribuite su 6 classi di componenti metallici: bracket black, bracket brown, bracket white, connector, metal plate, tubes. Il dataset è suddiviso in:
\begin{itemize}
\item \textbf{Training set}: 888 campioni (solo normali)
\item \textbf{Test set}: 458 campioni (176 normali, 282 anomali)
\end{itemize}

Le anomalie coprono scenari tipici: scratches (graffi), holes (buchi), missing parts (parti mancanti), rust (ruggine). Ogni immagine anomala è annotata con una pixel-precise ground truth mask.

\paragraph{Metodi Valutati.} MPDD confrontò due categorie di metodi:
\begin{itemize}
\item \textbf{Feature extraction-based}: PatchCore \cite{roth2022towards}, CFLOW-AD \cite{gudovskiy2022cflow}, PaDiM \cite{defard2021padim}, Semi-Orthogonal \cite{liznerski2021explainable}, SPADE \cite{cohen2021sub}
\item \textbf{Reconstruction-based}: DAGAN \cite{di2021dagan}, Skip-GANomaly \cite{akcay2019skip}, ITAE \cite{you2022unified}
\end{itemize}


\paragraph{Gap di Performance.} 
I risultati mostrarono la superiorità degli approcci feature-based (PatchCore, PaDiM) rispetto ai reconstruction-based. Tuttavia, i metodi subivano un \textbf{calo di per\-for\-mance} su MPDD:
\begin{itemize}
    \item \textbf{Differenza AUROC image-level}: fino a \textbf{23.12\%} per singoli metodi
    \item \textbf{Differenza media}: \textbf{15.24\%} tra le performance aggregate sui due dataset
\end{itemize}

Questo gap dimostrò in modo empirico che la vera generalizzazione richiede test su dataset che riflettano le condizioni operative reali.

\paragraph{Conclusione e Impatto.} 
MPDD contribuì a spostare il focus della comunità dalla semplice ottimizzazione delle metriche su benchmark standardizzati verso la robustezza in condizioni reali. Il suo messaggio chiave è che per lo sviluppo e il miglioramento dei metodi di AD, è necessario testare questi metodi su dataset basati su applicazioni reali specifiche.

%%% VisA

\subsection{VisA / SPot-the-Difference (2022)}
\label{sec:visa}
VisA (Visual Anomaly) \cite{zou2022spot} introdusse sfide significative e scenari più complessi: oggetti con strutture dettagliate, istanze multiple nello stesso frame, e anomalie sia superficiali che strutturali.

\paragraph{Caratteristiche.} Il dataset contiene \textbf{10,821 immagini RGB} (9,621 normali, 1,200 anomale) di 12 oggetti organizzati in 3 domini:
\begin{enumerate}
\item \textbf{Complex structures}: quattro tipi di printed circuit boards (PCB), con strutture complesse contententi transistors, capacitors, chips,
etc. 
\item \textbf{Multiple instances}: più oggetti nella stessa immagine (Macaroni1, Macaroni2, Capsule, Candele). Le istanze di Macaroni2 e Capsules sono molto diverse in posizione e orientamento degli oggetti.
\item \textbf{Single instance}: singolo oggetto per immagine (Cashew, Chewing gum, Fryum e Pipe fryum)
\end{enumerate}

Le anomalie, generate manualmente per produrre risultati realistici,  coprono due categorie principali:
\begin{itemize}
\item \textbf{Surface defects}: scratches, dents, color spots, cracks
\item \textbf{Structural defects}: misplacement (componenti disallineati), missing parts (parti mancanti)
\end{itemize}

Ci sono 5-20 immagini per tipo di anomalia e ogni immagine può contenerne più di una. Tutte le anomalie sono annotate con image-level e pixel-level labels.

\paragraph{Benchmark Protocols.} VisA introdusse protocolli few-shot e high-shot in aggiunta al paradigma one-class tradizionale:
\begin{itemize}
    \item \textbf{1-class vs 2-class training}: 
        \begin{itemize}
            \item \textbf{1-class}: addestramento su una singola classe (paradigma standard). Il 90\% delle immagini normali viene assegnato come \textit{train set} mentre il restante 10\% e le immagini anomale costituiscono il \textit{test set}.
            \item \textbf{2-class}: addestramento su dati normali \textbf{e} anomali per valutare la \textit{cross-category generalization}
        \end{itemize}
    \item \textbf{High-shot vs Low-shot} (definiti da percentuali di split):
        \begin{itemize}
            \item \textbf{High-shot}: 60\% dei dati normali/anomali per training, 40\% per test
            \item \textbf{Low-shot}: 20\% dei dati normali/anomali per training, 80\% per test
        \end{itemize}
    \item \textbf{k-shot (k=5,10)} (campionamento assoluto dal Low-shot):
        \begin{itemize}
            \item Dal training set Low-shot, si campionano \textbf{k immagini} dalla classe normale e \textbf{k immagini} dalla classe anomala
            \item Totale training: \textbf{2k immagini} (es: 10 immagini totali per 5-shot)
            \item Le performance sono mediate su 5 run random per robustezza statistica
        \end{itemize}
\end{itemize}

\paragraph{SPot-the-Difference (SPD).} Il paper introdusse SPD, un metodo di self-supervised pre-training progettato per aumentare la sensibilità locale di metodi SSL (Self-Supervised Learning) come SimCLR \cite{chen2020simple}, MoCo \cite{he2020momentum}, SimSiam \cite{chen2021exploring}. 

SPD utilizza \textbf{SmoothBlend}, una tecnica di local augmentation che genera pseudo-anomalie mescolando patch da immagini diverse con smoothing ai bordi per creare transizioni realistiche. Durante il pre-training, il modello impara a rilevare queste perturbazioni locali, sviluppando una sensibilità utile per l'anomaly detection.

\paragraph{Metodi Valutati.} Confronto con patch-based methods come PatchCore \cite{roth2022towards} e PaDiM \cite{defard2021padim}, dimostrando che SPD migliora le performance nei regimi few-shot.

\paragraph{Metriche.} Oltre all'AUROC standard, VisA utilizzò \textbf{AUPR (Area Under Precision-Recall curve)}, utile in scenari sbilanciati dove le anomalie sono rare.

\paragraph{Contributo.} VisA sottolineò l'importanza di:
\begin{enumerate}
\item \textbf{Scenari multi-istanza}: più comuni rispetto a single-object
\item \textbf{Efficienza dati}: few-shot learning per ridurre i costi di acquisizione
\item \textbf{Pre-training self-supervised}: come alternativa al fine-tuning supervisionato
\end{enumerate}

% LOCO
\subsection{MVTec LOCO AD (2022)}
\label{sec:mvtec_loco}
MVTec Logical Constraints Anomaly Detection (MVTec LOCO AD) \cite{bergmann2022beyond} introdusse per la prima volta le \textbf{anomalie logiche} oltre a quelle strutturali.

\paragraph{Motivazione.} In molti scenari industriali, un prodotto può essere strutturalmente integro ma logicamente errato. Esempi tipici possono includere:
\begin{itemize}
\item \textbf{Montaggio errato}: componenti montati nella posizione sbagliata, orientamento invertito
\item \textbf{Violazione di vincoli}: combinazioni di colori non conformi, numero errato di elementi
\item \textbf{Errori di assemblaggio}: parti mancanti in un insieme
\end{itemize}

Queste anomalie non sono rilevabili attraverso l'analisi di texture o geometria locale ma richiedono comprensione contestuale e reasoning sulle relazioni spaziali tra componenti, andando oltre il rilevamento di difetti superficiali.

\paragraph{Caratteristiche.} Il dataset contiene immagini RGB di \textbf{5 categorie di og\-get\-ti} (Breakfast Box, Screw Bag, Pushpins, Splicing Connectors, Juice Bottle), ciascuna con vincoli logici specifici. Il dataset è suddiviso in:
\begin{itemize}
\item \textbf{Training set}: 1,772 immagini (anomaly-free)
\item \textbf{Validation set}: 304 immagini (anomaly-free)
\item \textbf{Test set}: 1,568 immagini (con e senza anomalie)
\end{itemize}

Le anomalie sono sia \textbf{strutturali} (difetti fisici) che \textbf{logiche} (violazioni di vincoli logici). Entrambi i tipi hanno annotazioni pixel-precise.

\paragraph{Metodi Valutati.} Benchmark di diversi approcci:
\begin{itemize}
\item \textbf{Autoencoders}: Deterministic (AE), Variational (VAE) \cite{kingma2014auto}, Memory-guided (MNAD) \cite{park2020learning}
\item \textbf{GAN-based}: f-AnoGAN \cite{schlegl2019f}
\item \textbf{Feature-based}: SPADE \cite{cohen2021sub}, Student-Teacher \cite{bergmann2020uninformed}
\item \textbf{GCAD (Global Context Anomaly Detection)}: metodo specificamente progettato per anomalie logiche
\end{itemize}

\paragraph{Metriche.}  
Oltre all'AUROC, LOCO AD introdusse \textbf{sPRO (saturated Per-Region Overlap)}, una generalizzazione della metrica PRO. 
Mentre PRO standard calcola la media dell'overlap (percentuale di pixel correttamente predetti) per ciascuna regione anomala, pesando tutte le regioni allo stesso modo indipendentemente dalla loro dimensione, \textbf{sPRO} satura il contributo di una regione una volta che l'overlap supera una certa soglia. 

Questo approccio è particolarmente adatto per anomalie logiche: ad esempio, per un componente mancante (come una puntina da disegno assente in un contenitore), è sufficiente che il metodo segmenti un'area delle dimensioni di una puntina all'interno dello scomparto vuoto, senza bisogno di segmentare perfettamente tutti i pixel della regione annotata (che potrebbe essere più grande). 

\paragraph{Risultati Chiave.} I risultati rivelarono che:
\begin{enumerate}
\item Metodi reconstruction-based (AE, VAE, GAN) hanno performance significativamente peggiori su anomalie logiche rispetto a strutturali
\item Approcci feature-based con reti pre-addestrate generalizzano meglio
\item Il rilevamento di anomalie logiche rimane una sfida aperta, con AUROC tipicamente 5-10\% inferiore rispetto a MVTec AD standard
\end{enumerate}

\paragraph{Impatto.} MVTec LOCO AD evidenziò la necessità di metodi con capacità di reasoning contestuale, spingendo lo sviluppo di approcci basati su attention mechanisms e vision-language models capaci di comprendere relazioni semantiche.

% Real IAD

\subsection{Real-IAD (2024)}
\label{sec:real_iad}
Real-world Industrial Anomaly Detection (Real-IAD) \cite{ding2024realiad} rappresenta un passo significativo per la riduzione del gap tra performance su benchmark controllati e scenari di deployment reale.

\paragraph{Caratteristiche.} Real-IAD è uno dei dataset più grandi disponibili:
\begin{itemize}
    \item \textbf{150,000 immagini ad alta risoluzione}
    \item \textbf{30 categorie} diverse di oggetti di diversi materiali (metallo, plastica, legno, ceramica e mixed materials)
    \item \textbf{Acquisizione multi-view}: ogni oggetto viene catturato da 5 angolazioni distinte (dall'alto e da quattro angoli simmetrici a 45 gradi)
\end{itemize}

Questa struttura multi-view consente di valutare non solo la capacità di rilevamento ma anche la robustezza dei metodi rispetto a variazioni di prospettiva.

\paragraph{Protocolli di Valutazione.} Real-IAD introduce due setting sperimentali:

\subparagraph{1. Unsupervised IAD (UIAD).} L'algoritmo impara pattern e strutture \textit{solo} da campioni normali (paradigma standard). Il training set contiene solo imamgini normali, mentreil testing set ha sia campioni normali che anomali. Benchmark:
\begin{itemize}
\item \textbf{Embedding-based}: PatchCore \cite{roth2022towards}, PaDiM \cite{defard2021padim}, CFlow \cite{gudovskiy2022cflow}
\item \textbf{Data-augmentation-based}: SimpleNet \cite{liu2023simplenet}, DeSTSeg \cite{zhang2023destseg}
\item \textbf{Reconstruction-based}: RD (Reverse Distillation) \cite{deng2022anomaly}, UniAD \cite{you2022unified}
\end{itemize}

\subparagraph{2. Fully Unsupervised IAD (FUIAD).} Estensione realistica del setting UIAD dove il training set contiene \textit{rumore} (una frazione sconosciuta di campioni anomali mescolati ai normali), riflettendo scenari dove non è possibile garantire purezza assoluta del training set.

FUIAD è particolarmente rilevante per deployment industriale, dove la certificazione manuale di ogni immagine di training può essere proibitiva. I metodi devono essere robusti a contaminazione parziale del training set.

\paragraph{Metriche.} \textbf{AUROC} (image-level e pixel-level) e \textbf{AUPRO (Area Under PRO curve)}, che integra la metrica PRO su tutti i threshold.

\paragraph{Risultati.}
\begin{enumerate}
\item PatchCore mantiene performance elevate ($>95\%$ AUROC) 
\item Metodi data-augmentation-based (SimpleNet) mostrano robustezza superiore in FUIAD
\item Multi-view fusion migliora significativamente la detection rispetto a single-view
\end{enumerate}

\paragraph{Impatto.} Real-IAD stabilì un nuovo standard di scala e realismo, spingendo la ricerca verso soluzioni che non solo massimizzano le metriche su benchmark controllati, ma mantengono robustezza e affidabilità in condizioni operative realistiche. 

% mvtec 2
\subsection{MVTec AD 2 (2025)}
\label{sec:mvtec_ad2}
MVTec Anomaly Detection 2 (MVTec AD 2) \cite{bergmann2025mvtecad2} rappresenta l'evoluzione del benchmark MVTec AD originale, introducendo sfide che riflettono problemi aperti nella pratica industriale.

\paragraph{Composizione.} Il dataset contiene \textbf{8,004 immagini RGB ad alta ri\-so\-lu\-zio\-ne} distribuite su 8 scenari challenging:
\begin{enumerate}
\item \textbf{Bulk goods}: materiali sfusi (Wall
Plugs, Walnuts, Rice) con alta variabilità
\item \textbf{Textured objects}: oggetti con texture complesse e non uniformi
\item \textbf{Reflective metal objects}: superfici metalliche lucide con riflessi speculari
\item \textbf{Transparent objects}: oggetti trasparenti (vetro, plastica) difficili da ispezionare
\end{enumerate}

Ogni scenario presenta sfide specifiche che causano fallimenti nei metodi standard.

\paragraph{Caratteristiche Distintive.} MVTec AD 2 introduce innovazioni rispetto al predecessore:

\subparagraph{1. Anomalie ai Bordi dell'Immagine.} Le anomalie sono distribuite \textit{sull'intera altezza e larghezza} delle immagini, inclusi i bordi. Molti metodi (es. PatchCore con padding) hanno performance degradate vicino ai bordi.

\subparagraph{2. Condizioni di Illuminazione Realistiche.} Include acquisizioni con setup comuni nell'ispezione industriale:
\begin{itemize}
\item \textbf{Back-light illumination} (controluce): illuminazione da dietro l'oggetto
\item \textbf{Dark-field illumination} (campo scuro): illuminazione angolata che enfatizza difetti superficiali
\end{itemize}

Queste condizioni creano forte variabilità nell'aspetto di oggetti normali, rendendo più difficile distinguere anomalie da variazioni di illuminazione.

\subparagraph{3. Alta Variabilità Intra-Classe.} I dati normali presentano variazioni significative (orientamento, posizione, condizioni di luce), testando la robustezza dei modelli a distribution shift.

\paragraph{Metodi Benchmarkati.} MVTec AD 2 valutò metodi state-of-the-art rappresentativi di diverse famiglie:
\begin{itemize}
\item \textbf{Student-teacher}: EfficientAD \cite{batzner2024efficientad}, Reverse Distillation \cite{deng2022anomaly}, RD++ \cite{tien2023revisiting}
\item \textbf{Memory-bank}: PatchCore \cite{roth2022towards}
\item \textbf{Normalizing Flow}: MSFlow \cite{yu2021fastflow}
\item \textbf{Synthetic anomalies}: SimpleNet \cite{liu2023simplenet}, Dual Subspace Re-Projection
\end{itemize}

\paragraph{Nuove Metriche.} MVTec AD 2 introdusse metriche più rigorose:
\begin{itemize}
\item \textbf{AU-PRO (FPR 5\%)}: threshold-independent metric che integra PRO fino a un False Positive Rate del 5\%, enfatizzando precisione in regimi operativi realistici
\item \textbf{F1-score}: metrica threshold-dependent che bilancia precision e recall
\end{itemize}

\paragraph{Risultati Chiave.} I risultati rivelarono:
\begin{enumerate}
\item Nessun metodo raggiunge $>90\%$ AU-PRO su tutte le categorie (vs $>98\%$ AUROC su MVTec AD originale)
\item Oggetti trasparenti e riflettenti causano le maggiori difficoltà (AU-PRO $<70\%$)
\item Condizioni di illuminazione variabili aumentano significativamente i falsi positivi
\item EfficientAD e RD++ mostrano migliore robustezza complessiva
\end{enumerate}

\paragraph{Contributo.} MVTec AD 2 evidenziò che:
\begin{itemize}
\item Performance elevate su MVTec AD non garantiscono successo in scenari reali complessi
\item Sono necessari metodi robusti a distribution shift e condizioni di acquisizione variabili
\item Oggetti con proprietà ottiche speciali (trasparenza, riflessione) richiedono approcci dedicati
\end{itemize}
